\section{Introduction}
\label{sec:introduction}

Simulation is where underwater autonomy is developed, and the water is
where it is judged. Between the two sits a gap that the field usually
reports as a number -- a success rate that dropped, a tracking error that
grew -- and rarely as a mechanism. This paper reports the mechanism for
one camera-driven obstacle-avoidance stack, on one vehicle, in one pool,
and argues that the mechanism was not what the calibration effort had
been aimed at.

The study began from a conventional premise: if a simulator's mismatches
with reality are measured rather than guessed, closed-loop behaviour in
simulation should approach closed-loop behaviour in the water. We
measured three of them against the physical system. A static observation
model captured how often the real detector sees the obstacle at each
range and how badly its monocular range estimate is biased. A timing
model captured the rate at which observations arrive and, separately,
the burst structure of the gaps between them. A vehicle-response model
captured saturation, deadband, per-direction asymmetry and first-order
lag, applied downstream of the planner so that no calibration level ever
edits the controller. Stacking these on the historical simulator gives
four levels, S0 to S3, each strictly more realistic than the last.

We then did what a study of predictive validity requires and what small
robotics campaigns seldom do: we ran the simulated experiment first and
froze it. Eighty closed-loop runs -- four levels, two planners, two
obstacle geometries, five repetitions -- were executed, verified, hashed
and committed before the vehicle entered the water, so that the physical
campaign could only confirm or refute, never inform, the predictions.

The stack did not work on the vehicle. It detected the anchor in most
frames and commanded nothing. Making it act required seven changes made
over a single morning, and the pool was available for that one day. Nine
runs were flown with the committed planner and eight are admissible; the
dynamic-window baseline was never flown at all, because it requires a
pose estimate the vehicle carries no sensor to provide. The pool was
then drained. The physical dataset is closed and cannot be extended.

That outcome changes what the paper can honestly claim, and we let it.
A campaign that measures a different stack from the one that was frozen
cannot attribute a physical result to a calibration level, so we make no
claim of causal simulator predictivity and no claim about which planner
is better in the water. What the campaign can support is different, and
in our view more useful: an enumerated, file-level account of what had to
change, and a diagnosis of which changes mattered.

Two of the findings are geometric rather than statistical, and neither
was among the quantities we had spent months calibrating. First, the
detector that made the real system work is a fallback that accepts a
candidate only when its image height lies between $0.15$ and $0.90$ of
the frame; through the shared monocular estimator this is an observable
range window of \SI{0.29}{\metre} to \SI{2.11}{\metre}. The vehicle
started \SI{1.86}{\metre} from the anchor and its planner engaged at
\SI{1.80}{\metre}. The physical experiment therefore had no approach
phase: the planner was inside engagement range at the instant of
release, while its simulated counterpart approaches from
\SI{3.5}{\metre} and engages at \SI{1.5}{\metre}. Second, the planner
asks for a manoeuvre that the basin cannot contain -- a
\SI{2.5}{\metre} lateral offset and a \SI{4}{\metre} run past the
obstacle, in a pool \SIrange{2.7}{3.0}{\metre} wide. Nobody changed
those two parameters on the day, because nobody had noticed them.

A third finding concerns what a camera-only field campaign can measure at
all. No external tracking was recorded during the runs, so the vehicle's
trajectory, path length, lateral excursion and true clearance do not
exist for any physical run and cannot be recovered. We could have
integrated the commanded velocity and drawn plausible trajectories; we
do not, anywhere, because integrating an open-loop model of a
saturating, deadbanded vehicle in moving water manufactures precisely
the quantities under test. What remains is the command stream for all
eight runs, and -- recovered offline from the detector's own annotated
video -- the observation stream that reached the planner, including the
range the planner believed. The comparable surface between the two
domains is therefore the command domain, and we report it as such.

Finally, we audited our own frozen campaign and report what the audit
found, including a relay that converts image height to range at a
different field of view from the planner that inverts it, and a random
seed that is silently discarded. A frozen result whose defects are known
privately is worse than one whose defects are published.

Our contributions are:

\begin{enumerate}
\item a measurement-grounded calibration ladder for a camera-driven
  underwater avoidance loop, with every level fitted to paired
  observations of the physical vehicle rather than assumed;
\item a pre-registered, hashed simulated campaign executed before any
  physical run, together with a candid audit of its own defects;
\item a file-level account of a deployment gap -- seven changes, each
  with what forced it -- and the finding that the dominant obstacles were
  the geometry of the basin and the range window of the detector rather
  than any calibrated quantity;
\item a reproducible offline recovery of the physical perception stream
  from annotated video, and an explicit statement of which physical
  quantities are and are not measurable without external ground truth;
\item a post-deployment ablation that separates the effect of each
  deployment change, which a single pool day could not.
\end{enumerate}

Throughout, evidence gathered before the physical campaign is marked
\frozenmark, evidence gathered after it and with its outcome known is
marked \diagmark, and physical measurements are marked \realmark. The
distinction is load-bearing: the post-deployment simulations were
designed by people who already knew how the physical runs had gone, and
no amount of care makes that the same kind of evidence as a prediction.
